\documentclass[10pt,a4paper]{amsart}
\usepackage{amsmath,amssymb,amsthm,graphicx}
\usepackage{amsmath}
\usepackage{amssymb}
\usepackage{hyperref}
\newtheorem{Theorem}{Theorem}[section]
\newtheorem{Corollary}[Theorem]{Corollary}
\newtheorem{Lemma}[Theorem]{Lemma}
\newtheorem{Proposition}[Theorem]{Proposition}
{\theoremstyle{definition}
\newtheorem{Note}[Theorem]{Note}
\newtheorem{Definition}[Theorem]{Definition}
\newtheorem{Problem}[Theorem]{Problem}
}
\usepackage{cleveref}
\crefname{originallemma}{original lemma}{original lemmas}
\Crefname{originallemma}{Original lemma}{Original lemmas}
\makeatletter
\newlabel{lem:almostcom}{{1.25}{source}{}{sourcepointer-1}{}}
\newlabel{lem:almostcom@cref}{{[originallemma][0][]1.25}{[1][0][]source}{}{}{}}
\makeatother
\crefname{originalcorollary}{original corollary}{original corollarys}
\Crefname{originalcorollary}{Original corollary}{Original corollarys}
\makeatletter
\newlabel{lem:com}{{1.27}{source}{}{sourcepointer-2}{}}
\newlabel{lem:com@cref}{{[originalcorollary][0][]1.27}{[1][0][]source}{}{}{}}
\makeatother
\crefname{originallemma}{original lemma}{original lemmas}
\Crefname{originallemma}{Original lemma}{Original lemmas}
\makeatletter
\newlabel{lem:omegapower}{{1.32}{source}{}{sourcepointer-3}{}}
\newlabel{lem:omegapower@cref}{{[originallemma][0][]1.32}{[1][0][]source}{}{}{}}
\makeatother
\title{A basis for a complement of one filtration layer of the Askey-Wilson algebra in the next, built from Omega powers}
\author{Paul Terwilliger}
\date{doi:10.3842/SIGMA.2011.069; CC BY4.0}
\begin{document}
\maketitle
\noindent\textit{Editorial note.} All delivered mathematical text is the author's own, copied from the pinned publisher source sources/sigma-2011-069-askey-wilson/pinned-original-source.tex; no statement or proof step was rewritten, reordered or removed.\par\medskip
\noindent\textit{Editorial note.} Preamble provenance: the retained preamble is the source's own layout adapter plus the author's macro block at lines 10-18 of sources/sigma-2011-069-askey-wilson/pinned-original-source.tex (everything between the publisher class metadata and \textbackslash{}begin\{document\}, comments excluded). Nothing was added, no macro was invented, and the publisher document class is deliberately not reproduced; the excerpt serves only as preamble, never as a span coordinate.\par\medskip
\noindent\textit{Editorial note.} External references are navigation pointers only; every deliberately kept step of the author's argument stays inside the counted spans. The displayed numbers of those pointers are the author's own numbering under the author's \textbackslash{}newtheorem declarations (the source has no subsections, so numbering is section.index), computed mechanically by the lane's scanner.\par\medskip
\noindent\textit{Editorial note.} Three references to other labelled items occur inside the span, Corollary lem:com and Lemmas lem:omegapower and lem:almostcom; each is recorded once as a navigation stub with the matching kind and the author's own number, and every step that consumes them stays inside the counted span.\par\medskip
\noindent\textit{Editorial note.} The counted spans contain no bibliography citation, so no bibliography excerpt is attached to this card.\par\medskip
\noindent\textit{Editorial note.} Licence binding: the card's source\_version now carries the value recorded in this source's own licence evidence (sources/sigma-2011-069-askey-wilson/license-observation.json), namely 'SIGMA7(2011),069; publisher final form2011-07-09'. The adopt step of the publication line compares the two and refused the earlier revision; the card is aligned to the evidence rather than the evidence to the card.\par\medskip
\noindent\textit{Allowed standard prerequisites (editorial).}\begin{itemize}
\item A monomial basis for a complement of Delta\_\{n-1\} in Delta\_n. Same source, Corollary labelled lem:com; invoked inside the counted proof at lines 1676-1683. Supplies the six-tuple basis indexed by I\_n against which the Omega-power basis is compared.
\item Expressing a power of Omega in terms of monomials modulo the previous layer. Same source, Lemma labelled lem:omegapower; invoked inside the counted proof at lines 1696-1702. Replaces Omega\textasciicircum{}ell by q\textasciicircum{}\{ell\textasciicircum{}2\} A\textasciicircum{}ell B\textasciicircum{}ell C\textasciicircum{}ell modulo Delta\_\{n-1\}.
\item Moving the letters A, B, C past each other modulo the previous layer. Same source, Lemma labelled lem:almostcom; invoked inside the counted proof at lines 1703-1708. Converts A\textasciicircum{}ell B\textasciicircum{}ell C\textasciicircum{}ell into A\textasciicircum{}i B\textasciicircum{}j C\textasciicircum{}k form with the numerical factor q\textasciicircum{}\{2 j ell\}.
\end{itemize}\medskip
\ifdefined\hypertarget\hypertarget{sourcepointer-1}{}\fi
\noindent\textit{Source reference.} Original lemma 1.25: sources/sigma-2011-069-askey-wilson/pinned-original-source.tex, source lines 1703--1703. Navigation pointer for \textbackslash{}ref\{lem:almostcom\} used inside the counted span (line 1703); the argument that consumes it is retained in the span.\par\medskip
\ifdefined\hypertarget\hypertarget{sourcepointer-2}{}\fi
\noindent\textit{Source reference.} Original corollary 1.27: sources/sigma-2011-069-askey-wilson/pinned-original-source.tex, source lines 1676--1676. Navigation pointer for \textbackslash{}ref\{lem:com\} used inside the counted span (line 1676); the argument that consumes it is retained in the span.\par\medskip
\ifdefined\hypertarget\hypertarget{sourcepointer-3}{}\fi
\noindent\textit{Source reference.} Original lemma 1.32: sources/sigma-2011-069-askey-wilson/pinned-original-source.tex, source lines 1697--1697. Navigation pointer for \textbackslash{}ref\{lem:omegapower\} used inside the counted span (line 1697); the argument that consumes it is retained in the span.\par\medskip
\begin{Lemma}
\label{lem:basis2c}
For all integers $n\geq 1$
the following is a basis
for a complement of
$\Delta_{n-1}$ in $\Delta_n$:
\begin{gather*}
 A^iB^jC^k\Omega^\ell \alpha^r \beta^s \gamma^t, \qquad
i,j,k,\ell, r,s,t\geq 0,
\qquad ijk=0,
 \qquad
i+j+k+3\ell+r+s+t = n.
\end{gather*}
\end{Lemma}

\begin{proof}
Let $\mathbb I_n$ denote the set consisting of the 6-tuples
$(i,j,k,r,s,t)$ of nonnegative integers whose sum is~$n$.
By Corollary~\ref{lem:com} the
following
is a basis for a complement of
$\Delta_{n-1}$ in $\Delta_n$:
\begin{gather*}
A^iB^jC^k\alpha^r\beta^s\gamma^t, \qquad
(i,j,k,r,s,t) \in {\mathbb I}_n.
\end{gather*}
Let $\mathbb J_n$ denote the set consisting of the 7-tuples
$(i,j,k,\ell,r,s,t)$ of nonnegative integers such that
$ijk=0$ and
$i+j+k+3\ell+r+s+t = n$.
Observe that the map
\begin{gather*}
\mathbb J_n  \to  \mathbb I_n,
\qquad
(i,j,k,\ell,r,s,t)  \mapsto   (i+\ell,j+\ell,k+\ell,r,s,t)
\end{gather*}
is a bijection.
Suppose we are given $(i,j,k,\ell,r,s,t) \in \mathbb J_n$.
By
Lemma \ref{lem:omegapower},
$\Delta_{n-1}$ contains
\begin{gather*}
A^iB^jC^k\Omega^\ell \alpha^r \beta^s \gamma^t -
q^{\ell^2} A^iB^jC^k A^\ell B^\ell C^\ell \alpha^r \beta^s \gamma^t.
\end{gather*}
By Lemma~\ref{lem:almostcom},  $\Delta_{n-1}$ contains
\begin{gather*}
A^iB^jC^k A^\ell B^\ell C^\ell \alpha^r \beta^s \gamma^t
-
q^{2j \ell} A^{i+\ell}B^{j+\ell}C^{k+\ell}  \alpha^r \beta^s \gamma^t.
\end{gather*}
Therefore $\Delta_{n-1}$ contains
\begin{gather*}
A^iB^jC^k\Omega^\ell \alpha^r \beta^s \gamma^t -
q^{\ell(2j+ \ell)} A^{i+\ell}B^{j+\ell}C^{k+\ell}  \alpha^r \beta^s \gamma^t.
\end{gather*}
By these comments the following is a basis for a complement of
$\Delta_{n-1}$ in $\Delta_n$:
\begin{gather*}
A^iB^jC^k\Omega^\ell \alpha^r\beta^s\gamma^t, \qquad
(i,j,k,\ell,r,s,t) \in \mathbb J_n.
\end{gather*}
The result follows.
\end{proof}

\begin{thebibliography}{99}

\footnotesize\itemsep=0pt
\bibitem{alpern}
Alperin R.C.,
Notes: ${\rm  {PSL}}_2(\mathbb Z)= \mathbb Z_2 \star \mathbb Z_3$,
\href{http://dx.doi.org/10.2307/2324963}{{\it Amer. Math. Monthly}} {\bf 100} (1993), 385--386.
\end{thebibliography}
\end{document}
